\documentclass[11pt, twoside]{book}
\raggedbottom

\usepackage[utf8]{inputenc}

\usepackage[english]{babel}

\usepackage{microtype}
\usepackage{amsmath}
\usepackage{amssymb}
\usepackage{mathtools}
\usepackage{nicefrac}
\usepackage{pifont}

\def\spacebetweenparagraphs{2mm minus0.5mm}
\setlength{\parskip}{\spacebetweenparagraphs}

\def\horizontalindent{1.6em}
\setlength{\parindent}{\horizontalindent} % offset of the first line
\newlength\negparindent
\setlength{\negparindent}{-\parindent}

\renewcommand*\footnoterule{} % no line before footnotes

\usepackage{perpage}
\MakePerPage{footnote} % reset footnote numbering on every page

\usepackage[symbol,bottom,hang,flushmargin]{footmisc}

\setlength{\skip\footins}{2em plus1em} % space between the text body and the footnotes

\setlength{\footnotesep}{\baselineskip} % space between footnotes

\newcommand\oex{\hspace{1ex}} % horizontal space of the one ex
\newcommand\hex{\hspace{.5ex}} % one half of the ex
\newcommand\fex{\hspace{.2ex}} % one fifth of the ex
\newcommand\nfex{\hspace{-0.2ex}} % negative one fifth
\newcommand\thex{\hspace{.1ex}} % one tenth of the ex
\newcommand\nthex{\hspace{-0.1ex}} % negative one tenth

\makeatletter
\def\@fnsymbol#1{\ensuremath{\ifcase #1 %
	\or {*} %
	\or {*}\nfex{*} %
	\or {*}\nfex{*}\nfex{*} %
	\or {*}\nfex{*}\nfex{*}\nfex{*} %
	\else {\infty} %% \else \@ctrerr
\fi}}
\def\@makefnmark{\raisebox{.8ex}{\hbox{\normalfont\@thefnmark}}}

\renewcommand\footnotesize{%
   \@setfontsize\footnotesize\@ixpt{11}%
   \abovedisplayskip 8\p@ \@plus2\p@ \@minus4\p@
   \abovedisplayshortskip \z@ \@plus\p@
   \belowdisplayshortskip 4\p@ \@plus2\p@ \@minus2\p@
   \def\@listi{\leftmargin\leftmargini
               \topsep 4\p@ \@plus2\p@ \@minus2\p@
               \parsep 2\p@ \@plus\p@ \@minus\p@
               \itemsep \parsep}%
   \belowdisplayskip \abovedisplayskip
}
\makeatother

\let\originalfootnote\footnote
\renewcommand{\footnote}[1]{\originalfootnote{\hspace{.25em}\parbox[t]{.96\textwidth}{#1}}}

\renewcommand{\labelitemi}{\raisebox{-0.2ex}{\pmb{\checkmark}}}

\newcommand{\textsb}[1]{{\fontseries{sb}\selectfont #1}}

\newlength{\globalparindent}
\setlength{\globalparindent}{\parindent}

\newcommand{\indentpar}{\hspace*{\globalparindent}}

\newcommand{\lquote}[0]{“} % {``}
\newcommand{\rquote}[0]{”} % {''}
\newcommand{\inquotes}[1]{\lquote{#1}\rquote}

\newcommand*{\sprime}{^{\mkern1.5mu\prime}\mkern-2mu}
\newcommand*{\dprime}{^{\mkern1.5mu\prime\prime}\mkern-2mu}

\newcommand\domain{\mathsf{A}}
\newcommand\polynomialring[1]{\domain[#1]}
\newcommand\polynomialringintwo[2]{\domain[#1,#2]}
\newcommand\ring{\mathsf{R}}
\newcommand\cin{\nfex\in\nfex}
\newcommand\cnotin{\nfex\notin\nfex}
\newcommand\csubset{\nfex\subset\nfex}
\newcommand\ceq{\nfex=\nfex}
\newcommand\cneq{\nfex\neq\nfex}
\newcommand\cge{\nfex\ge\nfex}

\makeatletter
\newcommand{\raisemath}[1]{\mathpalette{\raisem@th{#1}}}
\newcommand{\raisem@th}[3]{\raisebox{#1}{$#2#3$}}
\makeatother

\newcommand\onehalf{\raisemath{.02em}{\frac{\raisemath{-0.16em}{1}}{2}}}
\newcommand\smalldisplaystyleonehalf{\scalebox{.8}{$\displaystyle \onehalf$}\thex}
\newcommand\smallerdisplaystyleonehalf{\scalebox{.7}{$\raisemath{.1em}{\displaystyle \onehalf}$}\thex}

%%%%%%\newcommand\imaginaryunit{\scalebox{.9}{$ \sqrt{\scalebox{1}[.98]{$\raisemath{-0.5em}{-1}$}} $}\thex}

\newcommand\imaginaryunit{%
  \raisebox{.1ex}{\scalebox{.9}{$%
    \hspace{-0.1em}\sqrt{\mathstrut\raisebox{-0.1ex}{$ \hspace{-0.1em} -1 $}}%
  $}\thex}%
}

\usepackage{tikz}

\newcommand\theproofbegins{\ensuremath{\tikz[baseline=-1ex] \draw [line width=.5pt, color=black, fill=white] (0, 0) circle (.8ex); \hspace{.6em}}}
\newcommand\theproofends{\ensuremath{\hspace{.6em} \tikz[baseline=-0.6ex] \draw [color=black, fill=black] (0, 0) circle (.8ex);}}

\usepackage{lettrine}

%%\input Nouveaud.fd
%%\newcommand*\nouveauinit{\usefont{U}{Nouveaud}{xl}{n}}
\input Zallman.fd
\newcommand*\zallmaninit{\usefont{U}{Zallman}{xl}{n}}

%%%%%%%%%%%%

\begin{document}

\thispagestyle{empty}

%
% intro
%

\noindent
\begin{center}
\scalebox{.95}{%
\begin{minipage}{\dimexpr(\linewidth)/95*100\relax}
\lettrine[lines=2, findent=.25ex]{\zallmaninit P}{olynomial}
\textls[-10]{%
division by a~linear polynomial lies at the~heart of~many results in~algebra.
Beginning with the~proof of the polynomial remainder theorem,
this note explores the~nature of roots, factors, and reducibility of~polynomials.
It then shows that the~ability to factor a~polynomial depends on the~choice of~co\-ef\-fi\-cient domain
and how passing to extension fields uncovers additional roots.
The~concluding example shows these ideas in~action and introduces the~concepts of splitting fields and root symmetries,
providing a~first glimpse of Galois theory.
}
\end{minipage}
}
\end{center}

\vspace{.25em}

%
% the polynomial remainder theorem
%

\noindent
\textsc{Theorem}.
{\itshape
In an integral domain~$\domain$,
for every polynomial~$P$ in the variable~$x$ with coefficients in~$\domain$, ${P \cin \polynomialring{x}}$,
and every ${r \cin \domain}$,
there exist unique ${\mathcal{Q} \cin \polynomialring{x}}$ %in the same polynomial ring
and~${\mathcal{R} \cin \domain}$
such that
}
\vspace{-0.1em}
\begin{equation*}
P(x) = (x - r) \fex \mathcal{Q}(x) + \mathcal{R}
\fex ,
\end{equation*}
%
\vspace{-2em}
\begin{flalign*}
\makebox[0pt][l]{\text{\itshape where}}
&&
\mathcal{R} = P(r)
\thex .
&&
\end{flalign*}

\noindent
\theproofbegins
\textsb{1. The existence}

\nopagebreak
Let ${P(x)}$ be a~polynomial in $\polynomialring{x}$ of degree~$n$,
%
\vspace{-0.5em}
\begin{equation*}
P(x) = a_n x^n + a_{n-1} x^{n-1} + \dots + a_1 x + a_0 = \scalebox{.88}{$ \displaystyle\sum_{i=0}^{n} $} \, a_i x^i
,
\vspace{-0.25em}
\end{equation*}
where each coefficient ${a_i \cin \domain}$. Evaluating it at ${x \ceq r}$ gives
%
\vspace{-0.25em}
\begin{equation*}
P(r) = a_n r^n + a_{n-1} r^{n-1} + \dots + a_1 r + a_0 = \scalebox{.88}{$ \displaystyle\sum_{i=0}^{n} $} \, a_i r^i
.
\vspace{-0.25em}
\end{equation*}
%
Consider the difference
%
\vspace{-0.5em}
\begin{equation*}
P(x) - P(r)
= \scalebox{.88}{$ \displaystyle\sum_{i=0}^{n} $} \, a_i x^i - \scalebox{.88}{$ \displaystyle\sum_{i=0}^{n} $} \, a_i r^i .
%%\vspace{-0.25em}
\end{equation*}
%
By the linearity of summation, the two expressions can be combined term-by-term.
Moreover, the constant terms for ${i \ceq 0}$ cancel out
(${x^0 \nfex = r^0 \nfex = 1}$ $\Rightarrow$
${a_0 x^0 \nfex - a_0 r^0 \ceq a_0 \nfex - a_0 \ceq 0}$).
The difference ${P(x) - P(r)}$ becomes
%
\vspace{-0.25em}
\begin{align*}
\scalebox{.88}{$ \displaystyle\sum_{i=0}^{n} $} \, a_i x^i - \scalebox{.88}{$ \displaystyle\sum_{i=0}^{n} $} \, a_i r^i
&= (a_0 \nfex - a_0) + \scalebox{.88}{$ \displaystyle\sum_{i=1}^{n} $} \Bigl( a_i x^i \nfex - a_i r^i \Bigr)
\\[-0.33em]
&= \scalebox{.88}{$ \displaystyle\sum_{i=1}^{n} $} \, a_i \bigl( x^i \nfex - r^i \bigr)
.
\vspace{-0.25em}
\end{align*}

\vspace{-0.75em}
Each term ${x^i \nfex - r^i}$ in the above sum is a~difference of two $i$-th powers.
The following lemma on the factorization of ${x^n \nfex - c^n}$ for any ${n \cge 1}$ shows that every such difference can be factored by ${(x - r)}$.
Applying the lemma term-by-term exposes a~common factor ${(x - r)}$ across all the terms, so it can be factored out of the~entire expression.

%%%%%%%%%%%%
\newpage
%%%%%%%%%%%%

\noindent
\textsc{Lemma}.
{\itshape
For any integer $n \cge 1$, the difference of two $n$-th powers, ${x^n \nfex - c^n\nfex}$,
can be factored by $(x - c)$.
That is, there exists a~polynomial ${S_{n-1}(x,c)}$
such that
}
\begin{equation*}
x^n \nfex - c^n = (x - c) \fex S_{n-1}(x,c)
\thex .
\end{equation*}

\noindent\theproofbegins
(by mathematical induction on $n$)

The \textsb{base case} ${n \ceq 1}$ holds since
%
\begin{equation*}
x^1 \nfex - c^1 = (x - c) \fex (1)
\fex ,
\end{equation*}
%
where the constant ${1 = S_0(x,c)}$ is a~polynomial as well.

\textsb{Inductive Hypothesis}.
Assume for some arbitrary positive integer~${k \cge 1}$
there exists a~polynomial $S_{k-1}(x,c)$ such that
%
\begin{equation*}
x^k \nfex - c^k = (x - c) \fex S_{k-1}(x,c)
\fex .
\end{equation*}

\textsb{Induction Step}.
To show that the lemma holds for~${k + 1}$ (assuming it holds for~$k$),
%(then it is true for all ${n \cge 1}$)
consider
%
\begin{equation*}
x^{k+1} \nfex - c^{k+1}
,
\end{equation*}
%
that can be rewritten by adding ${0 = - x \thex c^k + c^k x}$ and grouping as
%
\begin{equation*}
x^{k+1} \nfex - x \thex c^k \nfex + c^k x - c^{k+1} \nfex
= \bigl( x^{k+1} \nfex - x \thex c^k \bigr) + \bigl( c^k x - c^{k+1} \bigr)
\thex .
\end{equation*}
%
Factoring out the common components from each group gives
%
\begin{equation*}
x \bigl( x^k \nfex - c^k \bigr) + c^k(x - c)
\fex .
\end{equation*}
%
Applying the inductive hypothesis to substitute ${(x - c) \fex S_{k-1}(x,c)}$ for ${\bigl( x^k \nfex - c^k \bigr)}$ yields
%
\begin{equation*}
x (x - c) \fex S_{k-1}(x,c) + c^k(x - c)
\fex .
\end{equation*}
%
Factoring out the common multiplier $(x - c)$ results in
%
\begin{equation*}
(x - c) \Bigl( x \fex S_{k-1}(x,c) + c^k \Bigr)
,
\end{equation*}
%
where ${x \fex S_{k-1}(x,c) + c^k \nfex = S_k(x,c)}$ is a~polynomial
---
since ${S_{k-1}(x,c)}$ is a~polynomial by the inductive hypothesis,
and polynomial rings are closed under multiplication and addition.

This completes the induction step, proving the lemma for all ${n \cge 1}$.
\hfill
\theproofends

\vspace{.5em}

%
% the polynomials, explicitly
%
\scalebox{.8}{Explicitly, the~$S(x,c)$ polynomials are}

\nopagebreak\vspace{-0.5em}
\noindent
\scalebox{.8}{$
\begin{alignedat}{2}
& \oex 1 &{}= S_0 \\[-0.2em]
x(1) + c^1 &= x + c & {}= S_1 \\[-0.2em]
x (x + c) + c^2 &= x^2 + xc + c^2 & {}= S_2 \\[-0.2em]
x \bigl( x^2 + xc + c^2 \bigr) + c^3 &= x^3 + x^2c + xc^2 + c^3 & {}= S_3 \\[-0.2em]
x \bigl( x^3 + x^2c + xc^2 + c^3 \bigr) + c^4 &= x^4 + x^3c + x^2c^2 + xc^3 + c^4 & {}= S_4 \\[-0.4em]
& \cdots \\[-0.4em]
x \bigl( x^{n-2} + x^{n-3} c + \dots + x c^{n-3} + c^{n-2} \bigr) + c^{n-1}
&= x^{n-1} + x^{n-2} c + \dots + x c^{n-2} + c^{n-1} \nfex & {}= S_{n-1} \\[-0.2em]
x \bigl( x^{n-1} + x^{n-2} c + \dots + x c^{n-2} + c^{n-1} \bigr) + c^n
&= x^{n} + x^{n-1} c + \dots + x c^{n-1} + c^{n} & {}= S_n
\end{alignedat}
$}

\noindent
\scalebox{.8}{In more compact notation with the summation symbol}
%
\vspace{-0.8em}\begin{equation*}
\scalebox{.8}{$%
S_{n}(x,c)
= \fex x^{n} + x^{n-1} c + \dots + x c^{n-1} + c^{n}
= \scalebox{.88}{$ \displaystyle\sum_{i=0}^{n} $} \, x^{n-i} c^i
$}
\vspace{-1.75em}
\end{equation*}

\noindent
\scalebox{.8}{and}

\vspace{-0.25em}
\scalebox{.8}{$
\begin{aligned}
x \fex S_{n-1}(x,c) + c^n
& = x \, \scalebox{.88}{$ \displaystyle\sum_{i=0}^{n-1} $} \, x^{n-1-i} c^i + c^n
= \scalebox{.88}{$ \displaystyle\sum_{i=0}^{n-1} $} \, x^{n-i} c^i + c^n
= \scalebox{.88}{$ \displaystyle\sum_{i=0}^{n} $} \, x^{n-i} c^i
= S_{n}(x,c)
\fex .
\end{aligned}
$}

%
% the exact factorization
%
\vspace{-1.25em}
\begin{flalign*}
\makebox[0pt][l]{\scalebox{.8}{The exact factorization is}}
&&
\scalebox{.8}{$
x^n \nfex - c^n = (x - c) \fex \scalebox{.88}{$ \displaystyle\sum_{i=0}^{n-1} $} \, x^{n-1-i} c^i
.
$}
&&
\end{flalign*}

%
% the properties
%
\vspace{-0.75em}
\noindent
\scalebox{.8}{%
\parbox[t]{\dimexpr(\linewidth)/8*10\relax}{%
Several interesting properties characterize ${ S_{n}(x,c) = \textstyle\sum_{i=0}^{n} x^{n-i} c^i }$
as a~\emph{complete homogeneous symmetric polynomial} in two variables\::
}}
\vspace{-1.2em}
\begin{itemize}
\item
\scalebox{.8}{%
\parbox[t]{\dimexpr(\linewidth-\leftmargin)/8*10\relax}{%
it is \emph{symmetric}, since
${ S_{n}(x,c) = \textstyle\sum_{i=0}^{n} x^{n-i} c^i
= \textstyle\sum_{i=0}^{n} c^{n-i} x^i = S_{n}(c,x) }$,
}}
\vspace{-0.6em}
\item
\scalebox{.8}{%
\parbox[t]{\dimexpr(\linewidth-\leftmargin)/8*10\relax}{%
it is \emph{homogeneous} of~degree~$n$,
because the degree of each monomial ${x^{n-i} c^i}$ is exactly ${n \!-\! i + i = n}$
for all~$i$,
}}
\vspace{-0.4em}
\item
\scalebox{.8}{%
\parbox[t]{\dimexpr(\linewidth-\leftmargin)/8*10\relax}{%
it is \emph{complete}
in the~sense that it contains every possible \hbox{$n$-th} degree monomial in the variables~$x$ and~$c$
exactly once, each with a~coefficient of~$1$.
}}
\end{itemize}

%
% the dual recurrence
%
\vspace{-0.5em}
\noindent
\scalebox{.8}{%
\parbox[t]{\dimexpr(\linewidth)/8*10\relax}{%
The~symmetry of~$S_{n}(x,c)$ is already hinted~at in the~induction step.
The~proof inserts the~zero ${0 = -x \thex c^k + c^k x}$, but one may equally insert the~zero ${0 = -x^k c + c \thex x^k}$.
Then ${x^{k+1} \nfex - x^k c + c \thex x^k \nfex - c^{k+1} \nfex = x^k(x - c) + c \bigl( x^k - c^k \bigr)}$
and the resulting recurrence will be
%
\vspace{-0.66em}\begin{equation*}
S_k(x,c) = x^k \nfex + c \fex S_{k-1}(x,c)
\fex .
\vspace{-0.5em}
\end{equation*}
%
The~recurrences
${x^k \nfex + c \fex S_{k-1}(x,c)}$
and~${x \fex S_{k-1}(x,c) + c^k \nfex}$
are dual under the~swap ${x \nfex \leftrightarrow \nfex c \fex}$.
}}

%%%%%%%%%%%%

\noindent
\vspace{-0.5em}
\begin{flalign*}
\makebox[0pt][l]{\indentpar\text{Returning to}}
&&
P(x) - P(r)
= \scalebox{.88}{$ \displaystyle\sum_{i=1}^{n} $} \, a_i \bigl( x^i \nfex - r^i \bigr)
,
&&
\end{flalign*}
%
the lemma yields\footnote{%
From a~deeper algebraic viewpoint,
${S_{i-1}(x,c)}$ from the lemma belongs to the~polynomial ring $\polynomialringintwo{x}{c}$.
Replacing the~indeterminate~$c$ by the~element ${r \cin \domain}$ is the evaluation homomorphism ${\polynomialringintwo{x}{c} \to \polynomialring{x}}$.
Thus ${S_{i-1}(x,r)}$ is a~polynomial in~${\polynomialring{x}}$.
}
%
\begin{equation*}
x^i \nfex - r^i = (x - r) \fex S_{i-1}(x,r)
\tag*{\makebox[0pt][r]{for every ${i \cge 1 \fex}$.}}
\end{equation*}
%
Substituting these factorizations term-by-term gives
%
\begin{equation*}
P(x) - P(r)
= \scalebox{.88}{$ \displaystyle\sum_{i=1}^{n} $} \, a_i \bigl( x - r \bigr) S_{i-1}(x,r)
= \bigl( x - r \bigr) \scalebox{.88}{$ \displaystyle\sum_{i=1}^{n} $} \, a_i \fex S_{i-1}(x,r)
.
\end{equation*}

\vspace{-0.5em}
Since each $S_{i-1}(x,r)$ is a~polynomial in~${\polynomialring{x}}$ and polynomial sums are polynomials,
the~sum ${ \sum_{i=1}^{n} a_i \fex S_{i-1}(x,r) }$
belongs to the~same polynomial ring~$\polynomialring{x}$.
%
Define
%
\vspace{-0.1em}
\begin{equation*}
\scalebox{.88}{$ \displaystyle\sum_{i=1}^{n} $} \, a_i \fex S_{i-1}(x,r) = \mathcal{Q}(x) \in \polynomialring{x}
\fex .
\end{equation*}

\vspace{-1.75em}
\begin{flalign*}
\makebox[0pt][l]{\indentpar\text{Hence}}
&&
P(x) - P(r)
= \bigl( x - r \bigr) \mathcal{Q}(x)
&&
\\[.4em]
\makebox[0pt][l]{\text{or equivalently}}
&&
P(x) = \bigl( x - r \bigr) \mathcal{Q}(x) + P(r)
&&
\end{flalign*}

\vspace{-0.75em}
\noindent
--- the latter exhibits the intended representation with ${\mathcal{R} = P(r) \in \domain \fex}$.
Ergo, such~${\mathcal{Q} \cin \polynomialring{x}}$ and~${\mathcal{R} \cin \domain}$ exist.
\hfill
\theproofends

\vspace{.25em}
It remains to show that
${\mathcal{Q}(x)}$ and~${\mathcal{R}}$
are unique.

%%%%%%%%%%%%

\vspace{.5em}

\noindent
\theproofbegins
%
\textsb{2. The uniqueness}

\nopagebreak
Assume that there exist polynomials
${Q\sprime, Q\dprime \cin \polynomialring{x}}$
and constants
${R\sprime, R\dprime \cin \domain}$
such that
%
\vspace{-0.25em}
\begin{equation*}
P(x) = (x-r) \fex Q\sprime(x) + R\sprime
\quad \text{and} \quad
P(x) = (x-r) \fex Q\dprime(x) + R\dprime
.
\vspace{-0.2em}
\end{equation*}
%
Subtracting the~second representation from the~first gives
%
\vspace{-0.25em}
\begin{equation*}
(x-r) \fex Q\sprime(x) + R\sprime
- (x-r) \fex Q\dprime(x) - R\dprime
= 0
\fex ,
\vspace{-0.25em}
\end{equation*}
%
which simplifies to
%
\vspace{-0.5em}
\begin{equation*}
(x-r) \bigl( Q\sprime(x) \nfex - Q\dprime(x) \bigr) \nfex
= R\dprime \nfex - R\sprime
.
\end{equation*}

\vspace{-0.5em}
Evaluating this polynomial identity at~${x \ceq r}$,
the~factor~${(x - r)}$ becomes ${r - r = 0}$,
so the~left-hand side vanishes
and reveals\footnote{%
Since ${R\sprime, R\dprime \cin \domain}$,
the~difference~${R\dprime \nfex - R\sprime}$ is an~element of~$\domain$ itself\:---
equivalently, a~constant polynomial in~$\polynomialring{x}$.
Accordingly, once evaluation at~${x \ceq r}$ yields
${0 = R\dprime \nfex - R\sprime}$,
it follows that ${R\sprime \ceq R\dprime}$ as elements of~$\domain$.
In other words,
${R\sprime \ceq R\dprime}$ is an equality in~$\domain$ itself,
not merely equality obtained after evaluating the polynomial identity at~${x \ceq r}$.
}
that the constant difference is zero\::
%
\vspace{-0.25em}
\begin{equation*}
0 = R\dprime \nfex - R\sprime ,
\end{equation*}
%
\vspace{-2em}
\begin{flalign*}
\makebox[0pt][l]{\text{that is}}
&&
R\sprime = R\dprime .
&&
\end{flalign*}

\vspace{-0.5em}
Returning to the polynomial identity before evaluation
and putting ${R\sprime \ceq R\dprime}$ into it
yields
%
\vspace{-0.25em}
\begin{equation*}
(x - r) \bigl(Q\sprime(x) \nfex - Q\dprime(x)\bigr) \nfex = 0
\fex .
\end{equation*}

\vspace{-0.25em}
Because $\domain$ is an~integral domain,
the polynomial ring~$\polynomialring{x}$ is an~integral domain as well
and therefore has no zero divisors.
Moreover, ${x - r}$ is a~nonzero polynomial in~$\polynomialring{x}$.
Since its product with
${Q\sprime(x) \nfex - Q\dprime(x)}$
is the zero polynomial,
the absence of zero divisors in~$\polynomialring{x}$
implies that
%
\begin{equation*}
Q\sprime(x) \nfex - Q\dprime(x) \nthex = 0
\fex ,
\end{equation*}
%
\vspace{-1.75em}
\begin{flalign*}
\makebox[0pt][l]{\text{giving}}
&&
Q\sprime(x) \nthex = Q\dprime(x)
\thex .
&&
\end{flalign*}

\vspace{-1.75em}
\begin{flalign*}
\makebox[0pt][l]{\indentpar\text{Thus}}
&&
Q\sprime \ceq Q\dprime
\quad \text{and} \quad
R\sprime \ceq R\dprime
,
&&
\end{flalign*}
%
showing that the intended representation
%
\begin{equation*}
P(x) = (x - r) \fex \mathcal{Q}(x) + \mathcal{R}
\end{equation*}
%
is unique in the~choice of~$\mathcal{Q}$ and~$\mathcal{R}$.
\hfill
\theproofends

\vspace{1em}

%
% the definition of root
%

\noindent
\textsc{Definition}.
{\itshape
Let ${P \cin \polynomialring{x}}$ and ${r \cin \domain}$.
If ${P(r) \ceq 0}$, then ${r}$ is called a~\textbf{root} (or \textbf{zero}) of the polynomial~$P$.
Equivalently, ${r}$ is a~solution of the polynomial equation ${P(x) \nfex = 0}$.
}

\vspace{.5em}

%
% R=0 for a root
%

\noindent
\textsc{Corollary}.
{\itshape
If ${r \cin \domain}$ is a~root~(solution) of the polynomial equation ${P(x) \nfex = 0}$,
then there exists unique ${\mathcal{Q} \cin \polynomialring{x}}$ such that
}
\begin{equation*}
P(x) = (x-r) \mathcal{Q}(x)
\thex .
\end{equation*}

\vspace{-0.25em}
\noindent
\theproofbegins
By the previous theorem,
there exists unique ${\mathcal{Q} \cin \polynomialring{x}}$ such that
%
\begin{equation*}
P(x) = (x-r) \mathcal{Q}(x) + P(r)
\thex .
\end{equation*}
%
Since ${r}$ is a~root, ${P(r) \ceq 0}$.
Therefore
%
\begin{equation*}
P(x) = (x-r) \mathcal{Q}(x)
\thex .
\tag*{\theproofends}
\end{equation*}

%
% analogue with numerical division
%

In this case, $(x-r)$ is called a~\emph{factor} of~$P(x)$,
just as one says that a~number~$b$ is a~factor of another number~$a$ when~$b$ divides~$a$.
Hence every root~$r \in \domain$ gives a~linear factor~$(x-r)$ of~the~polynomial.

\vspace{.5em}

%
% mention that roots may appear only after enlarging the domain to a suitable extension field
%

%%\noindent
%%\textsc{Remark}.
The above definition concerns roots that belong to the same coefficient domain~$\domain$.
However, a~polynomial in $\domain[x]$ may fail to have roots inside~$\domain$
and acquire them only in a~larger domain or~field containing~$\domain$.
%%For instance, ${x^2 \nfex - 2 = 0}$ has no root in~$\mathbb{Q}$ but has the~roots ${\pm \sqrt{2}}$ in~$\mathbb{R}$.
Whenever a~root~$r$ exists in such an~extension, the~same theorem gives factorization
${P(x) = (x-r) \mathcal{Q}(x)}$
in the polynomial ring over that larger domain.

\vspace{.5em}

%
% reducibility / irreducibility
%

\noindent
\textsc{Definition}.
{\itshape
A~nonconstant polynomial ${P \cin \domain[x]}$ is called \textbf{reducible over~$\domain$}
if there exist nonconstant polynomials
${A(x), B(x) \cin \domain[x]}$
such that
}
\vspace{-0.25em}
\begin{equation*}
P(x) = A(x) \thex B(x)
\hspace{.1ex} .
\vspace{-0.25em}
\end{equation*}
{\itshape
If no such factorization exists, then ${P}$ is called \textbf{irreducible over~$\domain$}.
}

\vspace{.5em}

\noindent
\textsc{Remark}.
Reducibility is always relative to the chosen coefficient domain.
A~polynomial may be irreducible over one domain and reducible over a~larger one.

\noindent
\ding{118}\oex
For example, the polynomial ${x^2 \nfex - 2}$ is irreducible over~$\mathbb{Q}$.
Since it is quadratic, any nontrivial factorization would have to be into two linear factors.
But it has no rational root, hence no linear factor in~$\mathbb{Q}[x]$.
However, it has the~roots ${\pm \sqrt{2}}$ in the real numbers,
and over~$\mathbb{R}$ the same polynomial becomes reducible\footnote{%
Although ${x^2 \nfex - 2}$ splits into linear factors over the~reals,
the~field~$\mathbb{R}$ is not the~\emph{smallest} extension of~$\mathbb{Q}$ in which this factorization becomes possible.
A~smaller field already suffices, namely ${\mathbb{Q}\bigl(\sqrt{2}\bigr)}$.
This will~be explained after the concluding example.}
%
\vspace{-1em}
\begin{equation*}
x^2 \nfex - 2 \fex = \bigl(x - \sqrt{2}\bigr) \bigl(x + \sqrt{2}\bigr)
.
\vspace{-0.25em}
\end{equation*}
%
Thus, a~factorization that is impossible over one coefficient domain may become possible after passing to a~larger domain containing the~roots.

\noindent
\ding{118}\oex
In contrast, the~polynomial ${2x - 1}$ is irreducible over~$\mathbb{Z}$.
Yet it has the~rational root~${\nicefrac{1}{2}}$,
therefore over~$\mathbb{Q}$ it factors as
%
\vspace{-0.5em}
\begin{equation*}
2x - 1 = \Bigl( x - \smalldisplaystyleonehalf \Bigr) \nthex \bigl(2\bigr)
,
\vspace{-0.5em}
\end{equation*}
%
but the second factor is a~constant polynomial.
Over~$\mathbb{Q}$, ${2x - 1}$ remains irreducible,
because it cannot be written as a~product of two nonconstant polynomials.
Here, enlarging the coefficient domain uncovers the~root, but doesn’t produce a~nonconstant factorization, so the polynomial can’t become reducible.

Every nonconstant linear polynomial has degree~1 and is irreducible over any integral domain.
Thus, only polynomials of~degree two or~higher can become reducible by acquiring a~linear factor through one of their roots.

\vspace{.5em}

%
% reducibility when a root in the domain and the degree is at least two
%

\noindent
\textsc{Corollary}.
{\itshape
If a~polynomial ${P \cin \domain[x]}$ has a~root~${r \cin \domain}$ and the~degree of~$P$ is at~least~$2$,
then ${P}$ is reducible over~$\domain$.
}

\nopagebreak
\noindent
\theproofbegins\oex
If ${r \cin \domain}$ is a~root of ${P(x) \nfex = 0}$,
then by the first corollary there exists ${\mathcal{Q} \cin \domain[x]}$
such that
%
\begin{equation*}
P(x) = (x-r)\mathcal{Q}(x)
\hspace{.1ex} .
\end{equation*}
%
Since the~degree of~$P$ is at~least~$2$, both ${x-r}$ and the~polynomial~${\mathcal{Q}}$ are nonconstant.
That is, ${P}$ is a~product of two nonconstant polynomials in~$\domain[x]$,
therefore it is reducible over~$\domain$.
\hspace*{\fill}
\theproofends

\vspace{.5em}

%
% reducible doesn’t mean it has a root in the domain
%

\noindent
\textsc{Remark}.
The~existence of a~root in~$\domain$ implies reducibility over~$\domain$,
but reducibility does not necessarily mean that a~polynomial has a~root in~$\domain$.
For example, over~$\mathbb{R}$ the~polynomial
%
\vspace{-0.25em}
\begin{equation*}
x^4 \nfex + 4
=
\bigl(x^2 \nfex - 2x + 2\bigr)\bigl(x^2 \nfex + 2x + 2\bigr)
\vspace{-0.25em}
\end{equation*}
%
is reducible, although it has no real root.

%
% example
%

\vspace{.5em}

\noindent
\textsc{Example}
{\itshape
of how to factor a~polynomial into all of its linear factors via sequential domain extension.
}

\noindent
\scalebox{.95}{%
\begin{minipage}{\dimexpr(\linewidth)/95*100\relax}
\textls[-10]{%
This example explores how the~factorization and~reducibility of a~polynomial depend on the chosen coefficient domain.
As the~domain is enlarged, the~polynomial acquires new roots, and each newly obtained root translates into a~new linear factor.
}
\end{minipage}
}

\vspace{.5em}

Consider the polynomial
%
\vspace{-0.25em}
\begin{equation*}
P(x) = 2x^3 \nfex - x^2 \nfex - 4x + 2
\fex ,
\vspace{-0.25em}
\end{equation*}
%
whose coefficients all belong to~$\mathbb{Z}$.
Thus initially
${P \cin \mathbb{Z}[x]}$.

\vspace{.5em}

\noindent
\textsb{1. Over the integers $\mathbb{Z}$}

By grouping the terms and applying the distributive property,
%
\vspace{-0.2em}
\begin{equation*}
2x^3 \nfex - x^2 \nfex - 4x + 2
= 2x \bigl(x^2 \nfex - 2\bigr) - \bigl(x^2 \nfex - 2\bigr)
= (2x - 1) \bigl(x^2 \nfex - 2\bigr)
.
\vspace{-0.2em}
\end{equation*}
%
While the~factor ${(2x - 1)}$ suggests ${x \ceq \nicefrac{1}{2}}$ as a~candidate root
and ${(x^2 \nfex - 2)}$ becomes zero at~${x \ceq \pm \sqrt{2}}$, none of these values belong to~$\mathbb{Z}$.
Hence the~polynomial~$P(x)$ is reducible over~$\mathbb{Z}$ yet has no roots in~$\mathbb{Z}$.

\vspace{.5em}

\noindent
\textsb{2. Enlarging the domain from~$\mathbb{Z}$ to~$\mathbb{Q}$}

Now regard the same polynomial as an~element of~$\mathbb{Q}[x]$, ${\mathbb{Z} \csubset \mathbb{Q}}$.
Then
%
\vspace{-0.5em}
\begin{equation*}
P \bigl( \thex \nicefrac{1}{2} \thex \bigr) \nfex = 0
\fex ,
\vspace{-0.5em}
\end{equation*}
%
thus ${\nicefrac{1}{2}}$ is a~root in~$\mathbb{Q}$.
By the first corollary,
there exists a~polynomial ${Q \cin \mathbb{Q}[x]}$
such that
%
\nopagebreak\vspace{-0.25em}
\begin{equation*}
P(x) = \Bigl( x - \smalldisplaystyleonehalf \Bigr) \thex Q(x)
\thex .
\end{equation*}
%
Indeed,
%
\vspace{-0.75em}
\begin{equation*}
2x - 1 = \nfex \Bigl( x - \smalldisplaystyleonehalf \Bigr) \nthex \bigl(2\bigr)
\thex .
\end{equation*}

\vspace{-0.25em}
\noindent
\begin{minipage}{\linewidth}
Thereby, after extending the coefficient domain to~$\mathbb{Q}$,
the polynomial
%
\vspace{-0.25em}
\begin{equation*}
P(x) = (2x - 1)\bigl(x^2 \nfex - 2\bigr) = 2 \bigl(x - \nicefrac{1}{2}\bigr)\bigl(x^2 \nfex - 2\bigr)
,
\vspace{-0.25em}
\end{equation*}
%
remains reducible over~$\mathbb{Q}$.
%
The~first enlargement of the coefficient domain doesn’t change the~reducibility of~$P$, as it was already reducible over~$\mathbb{Z}$.
It merely replaces the linear factor ${(2x - 1)}$ by its monic associate ${\bigl(x - \nicefrac{1}{2}\bigr)}$,
while the quadratic factor remains irreducible over~$\mathbb{Q}$.
\end{minipage}

\vspace{.5em}

\noindent
\textsb{3. Enlarging the domain further from~$\mathbb{Q}$ to~$\mathbb{R}$}

In~$\mathbb{Q}$, the remaining factor ${(x^2 \nfex - 2)}$ still lacks roots because ${\pm\sqrt{2} \cnotin \mathbb{Q}}$.
However, in~$\mathbb{R}$ it splits as
%
\vspace{-0.25em}
\begin{equation*}
x^2 \nfex - 2 = \bigl(x - \sqrt{2}\fex\bigr)\bigl(x + \sqrt{2}\fex\bigr)
.
\vspace{-0.5em}
\end{equation*}
%
Therefore, in~$\mathbb{R}[x]$
%
\nopagebreak\vspace{-0.5em}
\begin{equation*}
P(x)
= 2 \Bigl(x - \smalldisplaystyleonehalf\Bigr) \nfex \Bigl(x - \sqrt{2}\fex\Bigr) \nfex \Bigl(x + \sqrt{2}\fex\Bigr)
.
\end{equation*}

\vspace{-0.5em}
After the~second enlargement of the coefficient domain, every irreducible factor is linear.
The~polynomial now splits completely into linear factors, and all of its roots become visible\::
%
\vspace{-0.25em}
\begin{equation*}
\smalldisplaystyleonehalf \fex, \quad \sqrt{2} \fex, \quad -\sqrt{2}
\fex .
\end{equation*}

%
% afterexample
%

\vspace{-0.5em}

\noindent
\scalebox{.8}{%
\begin{minipage}{\dimexpr(\linewidth)/8*10\relax}
\indentpar
The~second extension above, from~$\mathbb{Q}$ to~$\mathbb{R}$, is not the smallest possible one.
To factor the remaining quadratic term ${x^2 \nfex - 2}$, it already suffices to adjoin the single irrational number~${\sqrt{2}}$ to~$\mathbb{Q}$.
This produces the~field
%
\vspace{-0.75em}
\begin{equation*}
\mathbb{Q}\bigl(\sqrt{2}\bigr)
= \bigl\{\, a + b\sqrt{2} \oex \big|\; a,b \in \mathbb{Q} \,\bigr\}
\thex .
\vspace{-0.2em}
\end{equation*}
\end{minipage}
}

\noindent
\scalebox{.8}{%
\begin{minipage}{\dimexpr(\linewidth)/8*10\relax}
Every element of this field has the form ${a+b\sqrt{2}}$ with rational coefficients ${a,b}$.
It is an extension of~$\mathbb{Q}$ since it contains all the original rational numbers~$\mathbb{Q}$ when ${b \ceq 0}$.
Moreover, it is closed under the four arithmetic operations\:---
anyone can add, subtract, multiply and divide elements of~${\mathbb{Q}\bigl(\sqrt{2}\bigr)}$
and always get another number from~${\mathbb{Q}\bigl(\sqrt{2}\bigr)}$ that fits the~same ${a + b\sqrt{2}}$ pattern\::
%
\vspace{-0.5em}
\begin{equation*}
\begin{aligned}
\bigl(a + b\sqrt{2}\bigr) \nfex + \nfex \bigl(c + d\sqrt{2}\bigr)
&= (a+c) + (b+d)\sqrt{2}
\fex , \\%%[.2em]
\bigl(a + b\sqrt{2}\bigr) \nfex - \nfex \bigl(c + d\sqrt{2}\bigr)
&= (a-c) + (b-d)\sqrt{2}
\fex , \\%%[.2em]
\bigl(a + b\sqrt{2}\bigr) \nthex \bigl(c + d\sqrt{2}\bigr)
&= (ac + 2bd) + (ad+bc)\sqrt{2}
\fex ,
\end{aligned}
\vspace{-0.1em}
\end{equation*}
\end{minipage}
}

\noindent
\scalebox{.8}{%
\begin{minipage}{\dimexpr(\linewidth)/8*10\relax}
division works too, because the~denominator can always be \inquotes{rationalized} by multiplying by the~conjugate ${a - b\sqrt{2}}$
%
\vspace{-1em}
\begin{equation*}
\begin{aligned}
\scalebox{.93}{$ \displaystyle\frac{1}{a+b\sqrt{2}} $}
&=
\scalebox{.93}{$ \displaystyle\frac{a-b\sqrt{2}}{(a+b\sqrt{2})(a-b\sqrt{2})} $}
=
\scalebox{.93}{$ \displaystyle\frac{a-b\sqrt{2}}{a^2 \nfex - 2b^2} $}
\\
&=
\scalebox{.93}{$ \displaystyle\frac{a}{a^2 \nfex - 2b^2} $} - \scalebox{.9}{$ \displaystyle\frac{b}{a^2 \nfex - 2b^2}\sqrt{2} $}
\fex ,
\end{aligned}
\vspace{-0.5em}
\end{equation*}
%
provided ${a + b\sqrt{2} \neq 0}$, that is ${a}$ and~${b}$ are not both zero. %% ${a \cneq 0}$ or~${b \cneq 0}$
\end{minipage}
}

\noindent
\scalebox{.8}{%
\begin{minipage}{\dimexpr(\linewidth)/8*10\relax}
\indentpar
Since ${\mathbb{Q}\bigl(\sqrt{2}\bigr)}$ is a~field containing~$\mathbb{Q}$ and the roots ${\pm\sqrt{2}}$ of ${x^2 \nfex - 2}$,
the polynomial from the~example splits completely already in~${\mathbb{Q}\bigl(\sqrt{2}\bigr)[x]}$,
%
\vspace{-0.33em}
\begin{equation*}
P(x) \nthex = \thex 2 \Bigl(x - \smalldisplaystyleonehalf\Bigr) \nfex \Bigl(x - \sqrt{2}\fex\Bigr) \nfex \Bigl(x + \sqrt{2}\fex\Bigr)
.
\vspace{-0.1em}
\end{equation*}
\end{minipage}
}

\noindent
\scalebox{.8}{%
\begin{minipage}{\dimexpr(\linewidth)/8*10\relax}
\indentpar
The field ${\mathbb{Q}\bigl(\sqrt{2}\bigr)}$ also carries a~natural symmetry and can equally well be written as ${\mathbb{Q}\bigl(-\sqrt{2}\bigr)}$,
where ${\sqrt{2}}$ and ${-\sqrt{2}}$ can be swapped without changing the field itself.
Under this interchange, the two linear factors ${\bigl(x - \sqrt{2}\bigr)}$ and~${\bigl(x + \sqrt{2}\bigr)}$ are exchanged,
while their product
${\bigl(x + \sqrt{2}\bigr) \bigl(x - \sqrt{2}\bigr) \nfex = \thex x^2 \nfex - 2}$
remains the same.
Thus the two roots play completely symmetric roles inside the splitting field.
\end{minipage}
}

\noindent
\scalebox{.8}{%
\begin{minipage}{\dimexpr(\linewidth)/8*10\relax}
\indentpar
What emerges here is a~core idea of advanced algebra\:---
rather than enlarging the coefficient domain arbitrarily, one often seeks the \emph{smallest} field extension in which a~polynomial splits into linear factors.
Such a~minimal field, called a~\emph{splitting field}, contains exactly the~roots and the arithmetic combinations generated by them,
making the intrinsic symmetries among the roots visible without distraction of a~much larger ambient field.
In the present example, the~splitting field of~${(x^2 \nfex - 2)}$ over~$\mathbb{Q}$ is precisely ${\mathbb{Q}\bigl(\sqrt{2}\bigr)}$.
\end{minipage}
}

\noindent
\scalebox{.8}{%
\begin{minipage}{\dimexpr(\linewidth)/8*10\relax}
\indentpar
Another familiar example of such an~extension is the passage from~$\mathbb{R}$ to~$\mathbb{C}$.
The~polynomial ${x^2 \nfex + 1}$ has no root in~$\mathbb{R}$,
but adjoining a~root of it, $\imaginaryunit$, produces the larger field
${\mathbb{C} \ceq \mathbb{R}\bigl(\imaginaryunit\bigr)\nthex}$,
in which ${x^2 \nfex + 1}$ splits into linear factors,
%
\vspace{-0.67em}
\begin{equation*}
x^2 \nfex + 1 = \nfex \bigl( x - \imaginaryunit \fex \bigr) \bigl( x + \imaginaryunit \fex \bigr)
.
\end{equation*}
\end{minipage}
}

\noindent
\scalebox{.8}{%
\begin{minipage}{\dimexpr(\linewidth)/8*10\relax}
\indentpar
The systematic study of such splitting fields,
of the symmetries that permute polynomial roots while preserving the surrounding algebraic structure,
and of what those symmetries reveal about the polynomial itself,
leads to~\emph{Galois theory}.
\end{minipage}
}

\end{document}
